% Capítulo 5: Presentación de análisis, resultados y discusión

\chapter{Presentación de análisis, resultados y discusión}

\label{Chapter5}

\lhead{Capítulo 5. \emph{Resultados}}

\textit{[En un párrafo explica brevemente los contenidos de este capítulo antes de la distribución en distintos apartados. Se trata de la descripción de los resultados obtenidos y del trabajo de campo realizado. Se recomienda un tamaño de 5 a 15 páginas para este capítulo, pero depende del proyecto y de los datos obtenidos, por lo que la variabilidad posible es amplia.]}

%----------------------------------------------------------------------------------------

\section{Análisis descriptivo de la muestra}

\textit{[Describe la muestra seleccionada en función de las distintas variables consideradas y otros datos sociodemográficos de interés.]}

%----------------------------------------------------------------------------------------

\section{Resultados obtenidos}

\textit{[En este apartado se muestran los resultados obtenidos tras la investigación.]}

%----------------------------------------------------------------------------------------

\section{Discusión y análisis de los resultados}

\textit{[Interpretación del significado de los resultados obtenidos y análisis de su utilidad futura.]}
